\documentclass[10pt,a4paper]{amsart}
\usepackage[margin=19mm]{geometry}
\usepackage{amsmath,amssymb,mathtools}
\usepackage[scr=rsfs,cal=euler]{mathalfa}
\usepackage{xfrac}
\usepackage[unicode,hidelinks,bookmarks=false]{hyperref}
\newcommand{\ie}{i.e.\ }
\newcommand{\cf}{cf.\ }
\newcommand{\resp}{respectively\ }
\newcommand{\Subsection}{\subsection}
\providecommand{\nobreakdash}{\nobreak\hbox{-}}
\setlength{\emergencystretch}{2em}
\multlinegap0pt
\usepackage{bm}
\usepackage{bbm}
\usepackage{dsfont}
\usepackage{xspace}
\usepackage{cancel}
\usepackage{mathtools}
\usepackage{amsthm}
\makeatletter
\@ifundefined{definition}{\newtheorem{definition}{Definition}}{}
\@ifundefined{lemma}{\newtheorem{lemma}{Lemma}}{}
\@ifundefined{theorem}{\newtheorem{theorem}{Theorem}}{}
\makeatother
\title[Explicit construction of infinite families of strongly regul]{Explicit construction of infinite families of strongly regular digraphs with parameters $((v+(2^{n+1}-4)t)2^{n-1}, k+(2^n-2)t, t, \lambda, t)$}
\author{Viktor A. Byzov, Igor A. Pushkarev}
\date{Source: arXiv:2509.25425v1 (CC BY 4.0)}
\begin{document}
\maketitle

\noindent\emph{Source and licence.} Explicit construction of infinite families of strongly regular digraphs with parameters $((v+(2^{n+1}-4)t)2^{n-1}, k+(2^n-2)t, t, \lambda, t)$. Version arXiv:2509.25425v1 (CC BY 4.0), distributed under the Creative Commons Attribution 4.0 International licence, as recorded in the arXiv licence metadata for this exact version and shown on the abstract page https://arxiv.org/abs/2509.25425v1. Full rights notice: licenses/arxiv-2509.25425-CC-BY-4.0.txt.\par\smallskip

\noindent\textit{Editorial scope.} Counted unit: the Theorem whose source label is \texttt{th:main} in the article (source0.tex lines 313--337, source label \texttt{th:main}), delivered together with the article's own complete original proof of it (source0.tex lines 339--518, one \texttt{proof} environment of 180 lines). No mathematics is written, repaired or re-derived: statement and proof are the article's own text line for line. Required context retained uncounted: def:main2 (lines 99--110); eq:Pndef (lines 125--128); lemma:propPn (lines 132--139); lemma:Pn\_newdef (lines 233--240); lemma:summBnCn (lines 279--288); lemma:propsBnCnPn (lines 295--301). Every definition, display and label of those lines is retained unchanged. Omitted from this package and not redistributed: 1--98, 111--124, 129--131, 140--232, 241--278, 289--294, 302--312, 338--338, 519--599. Nothing inside a retained excerpt is omitted or altered: every retained body is delivered exactly as the article's lines are written, so no omission receipt of any kind accompanies this package. Citations: the retained excerpts carry no citation command at all, so no bibliography entry of the article is transcribed and no reference is redistributed. Reference adaptation: none was needed and none was made; every reference inside the delivered unit resolves inside it, so no unresolved reference or citation is produced. Printed slips kept literal: no printed word, symbol, hypothesis or number of the counted statement or of its proof was altered. Layout and class: the article's own class is \texttt{article} with packages including arxiv, inputenc, fontenc, hyperref, url, booktabs, amsfonts, nicefrac, microtype, lipsum, graphicx, doi, amssymb, amsthm; the wrapper is the lane's pinned \texttt{amsart} editorial wrapper at 10pt on A4, which restates the theorem-like environments the retained text needs and adds the article's own macro definitions that the retained text needs (none), with extra wrapper packages (bm, bbm, dsfont, xspace, cancel, mathtools). The delivered numbering is the unit's own. Assets: no retained span contains a figure environment, a graphics-inclusion command, a PDF-page inclusion, a table or a TikZ picture; no image file is copied into this package, the package declares no asset dependency and no image file is redistributed with it. Licence: version arXiv:2509.25425v1 of this article is distributed under the Creative Commons Attribution 4.0 International licence, as recorded in the arXiv licence metadata for this exact version and shown on the abstract page https://arxiv.org/abs/2509.25425v1. A full rights notice is delivered at licenses/arxiv-2509.25425-CC-BY-4.0.txt. Characters: every retained excerpt is pure ASCII, so the delivered engine, XeLaTeX with its default Latin Modern text font, reports no missing character.
\begin{definition}
	\label{def:main2}
	A strongly regular directed graph with parameter set ($v$, $k$, $t$, $\lambda$, $\mu$) is a digraph whose adjacency matrix $A$ satisfies the following conditions:
    \begin{equation}
	\label{eq:main1}
	A^2 = t I_v + \lambda A + \mu (J_v - I_v - A),
	\end{equation}
	\begin{equation}
	\label{eq:main2}
	A J_v = J_v A = k J_v.
	\end{equation}
\end{definition}

\begin{equation}
\label{eq:Pndef}
P_n = J_{2^n, 1} \otimes K_{2^n} \otimes J_{t, t\cdot 2^n}.
\end{equation}

\begin{lemma}
\label{lemma:propPn}
The matrix $P_n$, defined by formula~\eqref{eq:Pndef}, is a square matrix of order $t \cdot 4^n$ possessing the following properties:
\begin{enumerate}
	\item $P_n J_{t\cdot 4^n} = J_{t\cdot 4^n} P_n = t\cdot2^n J_{t\cdot 4^n}$;
	\item $P_n^2 = t J_{t\cdot 4^n}$.
\end{enumerate} 
\end{lemma}

\begin{lemma}
\label{lemma:Pn_newdef}
Let $P'_n$ be the sequence of matrices defined as follows:
\begin{equation}
P'_1 = J_{2, 1} \otimes K_2 \otimes J_{t, 2t}, \quad P'_n = J_{2^n, 1} \otimes K_2 \otimes \alpha_{2^{n-1}}(P'_{n-1}) \otimes J_{1, 2}.
\end{equation}
Then the sequence $P'_n$ coincides with the previously defined sequence of matrices $P_n$.
\end{lemma}

\begin{lemma}
\label{lemma:summBnCn}
The matrices $B_n$ and $C_n$ have the following properties:
\begin{enumerate}
	\item in each row of the matrix $B_n$ there are exactly $t\cdot 2^n$ ones;
	\item in each column of the matrix $B_n$ there are exactly $k + (2^n - 2)t$ ones;
	\item in each row of the matrix $C_n$ there are exactly $k + (2^n - 2)t$ ones;
	\item in each column of the matrix $C_n$ there are exactly $t\cdot 2^n$ ones.			
\end{enumerate}
\end{lemma}

\begin{lemma}$~$
\label{lemma:propsBnCnPn}
\begin{enumerate}
\item The matrices $B_n$ and $P_n$ have the following property: if any row of these matrices is partitioned into consecutive blocks of length $t\cdot 2^n$, then there is exactly one block consisting entirely of ones, and all the remaining blocks consist entirely of zeros.
\item The matrices $C_n$ and $P_n$ have the following property: if any column of these matrices is partitioned into consecutive blocks of length $t\cdot 2^n$, then each block contains exactly $t$ ones.
\end{enumerate}
\end{lemma}

\begin{theorem}
\label{th:main}
Let $A_n$ be a sequence of matrices in which $A_1$ is the adjacency matrix of the digraph $\text{dsrg}(v, k, t, \lambda, t)$, and for $n\geqslant 1$ it holds that
\begin{equation}
\label{eq:block_reprAnp1}
A_{n+1} = \left(
\begin{array}{cccc}
A_n & 0 & B_n & 0 \\
0 & A_n & 0 & B_n \\
0 & C_n & 0 & P_n \\
C_n & 0 & P_n & 0
\end{array}
\right)
= \left(
\begin{array}{cc}
I_2\otimes A_n & I_2\otimes B_n \\
K_2\otimes C_n & K_2\otimes P_n
\end{array}
\right)
,
\end{equation}
\noindent where $B_n$, $C_n$, and $P_n$ are the previously defined sequences (it is assumed that the matrices $B_1$ and $C_1$ exist).

Then the matrices $A_n$ are adjacency matrices of directed strongly regular digraphs with parameter sets~\mbox{$((v+(2^{n+1}-4)t)2^{n-1}, k+(2^n-2)t, t, \lambda, t)$}.
\end{theorem}

\begin{proof}
Introduce a notation for the order of the matrix $A_n$ (it is straightforward to show that it will indeed be this):
\begin{equation}
    v_n = (v+(2^{n+1}-4)t)2^{n-1}.
\end{equation}

It is necessary to prove two statements (see Definition~\ref{def:main2}):
\begin{enumerate}
	\item in each row and in each column of the matrix $A_n$ there are exactly $k+(2^n-2)t$ ones;
	\item $A_n^2 + sA_n = tJ_{v_n}$.
\end{enumerate}

Prove the first statement by induction on $n$. For the matrices $A_1$ and $A_2$ this fact holds. Assume that in each row and each column of $A_n$ there are exactly $k+(2^n-2)t$ ones. Then, using Lemma~\ref{lemma:propPn}, Lemma~\ref{lemma:summBnCn}, and formula~\eqref{eq:block_reprAnp1}, it follows that in each row and each column of $A_{n+1}$ there is also the required number of ones.

The proof of the second statement will also be carried out by induction on $n$. The base case was proved in the previous section by constructing the adjacency matrix of the digraph~$\text{dsrg}(2v+8t, k+2t, t, \lambda, t)$ in the block form given in formula~\eqref{eq:block_reprAnp1}.

Assume that $A_n^2 + sA_n = tJ_{v_n}$. Transform the expression $A_{n+1}^2 + sA_{n+1}$:
\begin{equation}
A_{n+1}^2 + sA_{n+1} = 
\left(
\begin{array}{c|c|c|c}
	A_n^2+sA_n & B_n C_n & A_n B_n + sB_n & B_n P_n \\ \hline
	B_n C_n & A_n^2+sA_n & B_n P_n & A_n B_n + sB_n \\ \hline
	P_n C_n & C_n A_n + sC_n & P_n^2 & C_n B_n + sP_n \\ \hline
	C_n A_n + sC_n & P_n C_n & C_n B_n + sP_n & P_n^2
\end{array}
\right).
\end{equation}

It is necessary to prove that all blocks of the resulting matrix are of the form $tJ$. The blocks on the main diagonal have the required form by the induction hypothesis and by Lemma~\ref{lemma:propPn}. Thus, it remains to prove the following six equalities:
\begin{align}
B_n C_n &= tJ_{v_n}; \label{eq:system1}\\
B_n P_n &= tJ_{v_n, t\cdot{4^n}}; \label{eq:system2}\\
P_n C_n &= tJ_{t\cdot{4^n}, v_n}; \label{eq:system3}\\
A_n B_n + sB_n &= tJ_{v_n, t\cdot{4^n}}; \label{eq:system4}\\
C_n A_n + sC_n &= tJ_{t\cdot{4^n}, v_n}; \label{eq:system5}\\
C_n B_n + sP_n &= tJ_{t\cdot{4^n}} \label{eq:system6}.
\end{align}

The validity of equalities~\eqref{eq:system1}--\eqref{eq:system3} follows automatically from Lemma~\ref{lemma:propsBnCnPn}. Based on the induction hypothesis, prove equalities~\eqref{eq:system4}--\eqref{eq:system6}. In the proof, Lemmas~\ref{lemma:propPn} and~\ref{lemma:Pn_newdef} will also be used.

Equality~\eqref{eq:system4}:
\begin{multline}
(A_n + s I_{v_n})B_n = \left(
\begin{array}{c|c}
I_2 \otimes (A_{n-1} + sI_{v_{n-1}}) & I_2 \otimes B_{n-1} \\ \hline
K_2 \otimes C_{n-1} & K_2 \otimes P_{n-1} + sI_{2t\cdot 4^{n-1}}
\end{array}	
\right)\cdot
\left(
\begin{array}{c}
K_2 \otimes B_{n-1} \otimes J_{1, 2} \\ \hline
I_2 \otimes P_{n-1} \otimes J_{1, 2}
\end{array}
\right)=
\\
=\left(
\begin{array}{c}
(I_2 \otimes (A_{n-1} + sI_{v_{n-1}}))\cdot (K_2 \otimes (B_{n-1} \otimes J_{1, 2})) + (I_2 \otimes B_{n-1})\cdot (I_2 \otimes (P_{n-1} \otimes J_{1, 2})) \\ \hline
(K_2 \otimes C_{n-1})\cdot (K_2 \otimes (B_{n-1} \otimes J_{1, 2})) + (K_2 \otimes P_{n-1} + sI_{2t\cdot 4^{n-1}})\cdot (I_2 \otimes (P_{n-1} \otimes J_{1, 2}))
\end{array}
\right)=
\\
=\left(
\begin{array}{c}
K_2\otimes ((A_{n-1} + sI_{v_{n-1}})\cdot (B_{n-1}\otimes J_{1, 2})) + I_2\otimes (B_{n-1}\cdot (P_{n-1}\otimes J_{1, 2})) \\ \hline
I_2\otimes (C_{n-1}\cdot (B_{n-1}\otimes J_{1,2})) + K_2\otimes (P_{n-1}\cdot (P_{n-1}\otimes J_{1, 2})) + sI_2\otimes P_{n-1}\otimes J_{1, 2} 
\end{array}
\right)=
\\
=\left(
\begin{array}{c}
K_2\otimes ((A_{n-1} + sI_{v_{n-1}})\cdot B_{n-1}) \otimes J_{1, 2} + I_2 \otimes (B_{n-1}\cdot P_{n-1}) \otimes J_{1, 2}\\ \hline
I_2\otimes(C_{n-1}B_{n-1}+sP_{n-1})\otimes J_{1, 2} + K_2\otimes (P_{n-1}\cdot P_{n-1})\otimes J_{1, 2}
\end{array}
\right)=
\\
=\left(
\begin{array}{c}
K_2\otimes tJ_{v_{n-1}, t\cdot 4^{n-1}} \otimes J_{1, 2} + I_2 \otimes tJ_{v_{n-1}, t\cdot 4^{n-1}} \otimes J_{1, 2}\\ \hline
I_2\otimes tJ_{t\cdot 4^{n-1}} \otimes J_{1, 2} + K_2\otimes tJ_{t\cdot 4^{n-1}} \otimes J_{1, 2}
\end{array}
\right)=
\left(
\begin{array}{c}
J_2\otimes tJ_{v_{n-1}, t\cdot 4^{n-1}} \otimes J_{1, 2}\\ \hline
J_2\otimes tJ_{t\cdot 4^{n-1}} \otimes J_{1, 2}
\end{array}
\right) =\\= tJ_{v_n, t\cdot{4^n}}.
\end{multline}

Equality~\eqref{eq:system5}:
\begin{multline}
C_n(A_n+sI_{v_n}) =
\left(
\begin{array}{c|c}
J_{2^n, 1} \otimes I_2 \otimes \alpha_{2^{n-1}}(C_{n-1}) & J_{2^n, 1} \otimes I_2 \otimes \alpha_{2^{n-1}}(P_{n-1})
\end{array}
\right)\times\\
\times
\left(
\begin{array}{c|c}
I_2 \otimes (A_{n-1} + sI_{v_{n-1}}) & I_2 \otimes B_{n-1} \\ \hline
K_2 \otimes C_{n-1} & K_2 \otimes P_{n-1} + sI_{2t\cdot 4^{n-1}}
\end{array}	
\right)=\\
=\left(
\begin{array}{c|c}
\text{\footnotesize $((J_{2^n, 1} \otimes I_2) \otimes \alpha_{2^{n-1}}(C_{n-1}))\cdot (I_2 \otimes (A_{n-1} + sI_{v_{n-1}}))+((J_{2^n, 1} \otimes I_2) \otimes \alpha_{2^{n-1}}(P_{n-1}))\cdot(K_2 \otimes C_{n-1})$} & ~
\end{array}
\right.
\\
\left.
\begin{array}{c|c}
~ & \text{\footnotesize $((J_{2^n, 1} \otimes I_2) \otimes \alpha_{2^{n-1}}(C_{n-1}))\cdot (I_2 \otimes B_{n-1}) + ((J_{2^n, 1} \otimes I_2) \otimes \alpha_{2^{n-1}}(P_{n-1}))\cdot(K_2 \otimes P_{n-1} + sI_{2t\cdot 4^{n-1}})$}
\end{array}
\right)=
\\
=\left(
\begin{array}{c|c}
\text{\footnotesize $((J_{2^n, 1} \otimes I_2)\cdot I_2) \otimes (\alpha_{2^{n-1}}(C_{n-1})\cdot (A_{n-1} + sI_{v_{n-1}}))+((J_{2^n, 1} \otimes I_2)\cdot K_2) \otimes (\alpha_{2^{n-1}}(P_{n-1})\cdot C_{n-1})$} & ~
\end{array}
\right.
\\
\begin{array}{c|c}
~ & \text{\footnotesize $((J_{2^n, 1} \otimes I_2)\cdot I_2) \otimes (\alpha_{2^{n-1}}(C_{n-1}) \cdot B_{n-1})+((J_{2^n, 1}\otimes I_2) \cdot K_2) \otimes (\alpha_{2^{n-1}}(P_{n-1})\cdot P_{n-1}) + $}
\end{array}
\\
\left.
\begin{array}{c}
\text{\footnotesize $ + sJ_{2^n, 1}\otimes I_2\otimes \alpha_{2^{n-1}}(P_{n-1})$}
\end{array}
\right)=
\\
=\left(
\begin{array}{c|c}
J_{2^n, 1} \otimes I_2 \otimes \alpha_{2^{n-1}}(tJ_{t\cdot 4^{n-1}, v_{n-1}}) + J_{2^n, 1} \otimes K_2 \otimes \alpha_{2^{n-1}}(tJ_{t\cdot 4^{n-1}, v_{n-1}}) & ~
\end{array}
\right.\\
\left.
\begin{array}{c|c}
~ & J_{2^n, 1} \otimes I_2 \otimes \alpha_{2^{n-1}}(tJ_{t\cdot 4^{n-1}}) + J_{2^n, 1} \otimes K_2 \otimes \alpha_{2^{n-1}}(tJ_{t\cdot 4^{n-1}})
\end{array}
\right)=\\
=\left(
\begin{array}{c|c}
J_{2^n, 1} \otimes J_2 \otimes \alpha_{2^{n-1}}(tJ_{t\cdot 4^{n-1}, v_{n-1}}) &
J_{2^n, 1} \otimes J_2 \otimes \alpha_{2^{n-1}}(tJ_{t\cdot 4^{n-1}})
\end{array}
\right)=tJ_{t\cdot{4^n}, v_n}.
\end{multline}

Equality~\eqref{eq:system6}:
\begin{multline}
C_n B_n + sP_n =
\left(
\begin{array}{c|c}
J_{2^n, 1} \otimes I_2 \otimes \alpha_{2^{n-1}}(C_{n-1}) & J_{2^n, 1} \otimes I_2 \otimes \alpha_{2^{n-1}}(P_{n-1})
\end{array}
\right)\times\\
\times
\left(
\begin{array}{c}
K_2 \otimes B_{n-1} \otimes J_{1, 2} \\ \hline
I_2 \otimes P_{n-1} \otimes J_{1, 2}
\end{array}	
\right) + sJ_{2^n, 1} \otimes K_2 \otimes \alpha_{2^{n-1}}(P_{n-1}) \otimes J_{1, 2}=\\
= ((J_{2^n, 1} \otimes I_2) \otimes \alpha_{2^{n-1}}(C_{n-1}))\cdot (K_2 \otimes (B_{n-1} \otimes J_{1, 2})) + 
\\
 + ((J_{2^n, 1} \otimes I_2) \otimes \alpha_{2^{n-1}}(P_{n-1}))\cdot (I_2 \otimes (P_{n-1} \otimes J_{1, 2})) + sJ_{2^n, 1} \otimes K_2 \otimes \alpha_{2^{n-1}}(P_{n-1}) \otimes J_{1, 2}
=\\
= ((J_{2^n, 1} \otimes I_2)\cdot K_2) \otimes (\alpha_{2^{n-1}}(C_{n-1}) \cdot (B_{n-1} \otimes J_{1, 2})) + 
\\
 + ((J_{2^n, 1}\otimes I_2)\cdot I_2)\otimes (\alpha_{2^{n-1}}(P_{n-1})\cdot (P_{n-1}\otimes J_{1,2})) + sJ_{2^n, 1} \otimes K_2 \otimes \alpha_{2^{n-1}}(P_{n-1}) \otimes J_{1, 2}
=\\
= J_{2^n, 1} \otimes K_2 \otimes \alpha_{2^{n-1}}(tJ_{t\cdot 4^{n-1}}) \otimes J_{1, 2} + J_{2^n, 1} \otimes I_2 \otimes \alpha_{2^{n-1}}(tJ_{t\cdot 4^{n-1}}) \otimes J_{1, 2}
=\\= J_{2^n, 1} \otimes J_2 \otimes \alpha_{2^{n-1}}(tJ_{t\cdot 4^{n-1}}) \otimes J_{1, 2} = tJ_{t\cdot 4^n}.
\end{multline}

\end{proof}

\end{document}
